\section{RAG: Retrieval-Augmented Generation}

\eng{Retrieval-Augmented Generation} (\eng{RAG}) הוא מנגנון שמאפשר לסוכן
לגשת לידע חיצוני מבלי לאמן מחדש את המודל \cite{lewis2020rag}. בהרצאה
הושווה ה-\eng{RAG} ל``דיסק קשיח'': מאגר גדול המסודר וקטורית לחיפוש סמנטי.

% RAG retrieval flow
\begin{figure}[H]
\centering
\begin{tikzpicture}[
  node distance=1.1cm and 1.4cm,
  block/.style={draw=primary, fill=lightgray, rounded corners=3pt,
    minimum width=2.6cm, minimum height=0.9cm, align=center, font=\footnotesize},
  arr/.style={-{Stealth[length=2mm]}, thick, accentblue}
]
  \node[block] (prompt) {שאילתת\\משתמש};
  \node[block, right=of prompt] (embed) {הטמעה\\וקטורית};
  \node[block, right=of embed] (search) {חיפוש\\סמנטי};
  \node[block, below=of search] (docs) {מסמכים\\רלוונטיים};
  \node[block, left=of docs] (ctx) {\eng{Context}\\\eng{Window}};
  \node[block, left=of ctx] (llm) {\eng{LLM}};
  \node[block, below=of llm] (answer) {תשובה};
  \draw[arr] (prompt) -- (embed);
  \draw[arr] (embed) -- (search);
  \draw[arr] (search) -- (docs);
  \draw[arr] (docs) -- (ctx);
  \draw[arr] (prompt.south) |- (ctx.west);
  \draw[arr] (ctx) -- (llm);
  \draw[arr] (llm) -- (answer);
\end{tikzpicture}
\caption{זרימת \eng{RAG}: מהפרומפט לשליפה וקטורית, הזרקה לחלון הקשר ותשובת המודל}
\label{fig:rag-flow}
\end{figure}


\subsection{עקרון הפעולה}

כאשר משתמש שואל שאלה:
\begin{enumerate}[rightmargin=2em]
  \item הפרומפט מומר לוקטור (הטמעה --- \eng{embedding}).
  \item מתבצע חיפוש במאגר הוקטורים באמצעות מרחק קוסינוס (או מטריקה דומה).
  \item המסמכים הקרובים ביותר סמנטית נשלפים מהמאגר.
  \item המסמכים מצורפים לחלון ההקשר יחד עם הפרומפט.
  \item ה-\eng{LLM} מייצר תשובה המבוססת על המידע שנשלף.
\end{enumerate}

דוגמה מההרצאה: שאלה על בירת ישראל $\rightarrow$ חיפוש בוויקיפדיה ובמסמכים
רלוונטיים $\rightarrow$ הזרקה ל-\eng{Context Window} $\rightarrow$ תשובה מדויקת.

\subsection{מדוע לא לאמן מחדש?}

אימון מחדש של \eng{LLM} על מאגר ארגוני הוא יקר, איטי ומסובך לתחזוקה.
\eng{RAG} מאפשר:
\begin{itemize}[rightmargin=2em]
  \item עדכון ידע בזמן אמת (הוספת מסמכים למאגר).
  \item ציטוט מקורות (איזה מסמך שימש בתשובה).
  \item שליטה בהרשאות (מאגרים נפרדים לפי מחלקה).
\end{itemize}

\subsection{דוגמת זרימה (פסאודו-קוד)}

\begin{lstlisting}[language=Python,caption={Pipeline בסיסי של RAG}]
def rag_answer(question, vector_db, llm):
    q_vec = embed(question)
    docs = vector_db.search(q_vec, top_k=5)
    context = "\n".join(docs)
    prompt = f"Context:\n{context}\n\nQuestion: {question}"
    return llm.generate(prompt)
\end{lstlisting}

\begin{warningbox}
\eng{RAG} אינו מבטיח אמיתות --- אם המאגר מכיל מידע שגוי, המודל עלול
להמחיש אותו בביטחון. נדרשת בקרת איכות על המקורות ואימות תשובות.
\end{warningbox}

\subsection{רכיבי מערכת RAG}

\agenttable{
\caption{רכיבי מערכת \eng{RAG} טיפוסית}
\label{tab:rag-components}
\begin{tabular}{@{}p{3cm}p{4.5cm}p{4.5cm}@{}}
\toprule
\textbf{רכיב} & \textbf{תפקיד} & \textbf{דוגמה} \\
\midrule
Ingestion & פיצול ואינדוקס מסמכים & \eng{chunk} של 512 טוקנים \\
Embedding & המרה לוקטור & \eng{text-embedding-3-small} \\
Vector DB & אחסון וחיפוש & \eng{Chroma}, \eng{Pinecone} \\
Retriever & בחירת קטעים & \eng{top-k}, סף דמיון \\
Generator & תשובה סופית & \eng{GPT-4}, \eng{Claude} \\
\bottomrule
\end{tabular}
}

\subsection{קשר לקורס}

במטלות קודמות בקורס (למשל תרגום רב-שלבי) נעשה שימוש בהשוואת וקטורים
למדידת סטייה סמנטית. \eng{RAG} משתמש באותו עקרון --- מרחק וקטורי ---
למטרה אחרת: שליפת ידע במקום הערכת דמיון בין תרגומים.

\subsection{מגבלות}

\begin{itemize}[rightmargin=2em]
  \item איכות התלויה בגודל ואיכות ה-\eng{chunk}.
  \item חיפוש שגוי אם השאלה מנוסחת רחוק מהמסמכים.
  \item עלות הטמעה ואחסון למאגרים גדולים.
  \item צורך באבטחת מידע --- דליפת מסמכים רגישים לפרומפט.
\end{itemize}

\begin{insightbox}
בפרק 13 נראה כיצד \eng{RAG} יכול לתמוך בחקירת \eng{Ransomware} על ידי
אינדוקס דוחות \eng{CTI}, הערות כופר ו-\eng{IOCs} היסטוריים.
\end{insightbox}
